\documentclass[11pt]{article}

\usepackage[margin=1in]{geometry}
\usepackage{amsmath,amssymb}
\usepackage{mathtools}
\usepackage{enumitem}

\title{Power-Spectrum Computation With and Without Periodic Boundary Conditions}
\author{}
\date{}

\begin{document}
\maketitle

\section{The object being analyzed}

In the full Planck component-separation phase, each minibatch object is a
5-channel map
\[
  X \in \mathbb{R}^{N_b\times C\times N\times N},
  \qquad C=5.
\]
For one nuisance realization $i$, the target channels are
\[
  X_{\mathrm{target}}^{(i)}
  =
  \begin{bmatrix}
    d_Q\\
    n_Q^{(i)}\\
    d_U\\
    n_U^{(i)}\\
    d_I
  \end{bmatrix},
\]
and the running channels are
\[
  X_{\mathrm{run}}^{(i)}(u)
  =
  \begin{bmatrix}
    u_Q+n_Q^{(i)}\\
    d_Q-u_Q\\
    u_U+n_U^{(i)}\\
    d_U-u_U\\
    d_I
  \end{bmatrix}.
\]

The power-spectrum term computes, for each minibatch element, channel pair,
and radial Fourier bin,
\[
  P(X)\in\mathbb{C}^{N_b\times C\times C\times B}.
\]
Here $B$ is the number of radial Fourier bins. The two $C$ indices correspond
to a channel pair $(c_1,c_2)$.

\section{Common Fourier-bin construction}

Let
\[
  \widehat X_{a,c}(k)
\]
be the 2D Fourier transform of minibatch element $a$ and channel $c$, where
$k=(k_x,k_y)$ is a discrete Fourier pixel.

For each radial bin $b$, the code constructs a smooth radial mask
\[
  M_b(k).
\]
This mask selects Fourier modes whose radius $|k|$ lies near bin $b$.

The first common intermediate is the binned Fourier map
\[
  \widehat X^{(b)}_{a,c}(k)
  =
  M_b(k)\widehat X_{a,c}(k).
\]
In tensor form this has shape
\[
  N_b\times C\times B\times N\times N.
\]

\section{\texttt{PBC=True}: periodic power-spectrum estimate}

When \texttt{PBC=True}, the map is treated as periodic. In that case, Fourier
modes are taken as the natural orthogonal basis of the domain, and the code
stays entirely in Fourier space.

For a channel pair $(c_1,c_2)$ and bin $b$, the code computes
\[
  P^{\mathrm{pbc}}_{a,c_1,c_2,b}
  =
  \frac{
    \sum_k
    M_b(k)\widehat X_{a,c_1}(k)
    \overline{\widehat X_{a,c_2}(k)}
  }{
    \sum_k M_b(k)
  }.
\]

Equivalently,
\[
  P^{\mathrm{pbc}}_{a,c_1,c_2,b}
  =
  \left\langle
    \widehat X^{(b)}_{a,c_1}(k)
    \overline{\widehat X_{a,c_2}(k)}
  \right\rangle_{k,b},
\]
where $\langle\cdot\rangle_{k,b}$ denotes an average over Fourier pixels
weighted by the radial bin mask.

\subsection*{Interpretation}

This branch asks: for the 5-channel object, how much Fourier power or
cross-power does each selected channel pair have in radial bin $b$?

For example,
\[
  P_{0,0,b}
\]
is the binned auto-power of the first channel, and
\[
  P_{0,4,b}
\]
is the binned cross-power between the first channel and the auxiliary
$d_I$ channel.

\section{\texttt{PBC=False}: non-periodic power-spectrum estimate}

When \texttt{PBC=False}, the map is not assumed to wrap around at the edges.
This matters because FFTs implicitly treat the input as periodic. If the left
and right edges, or top and bottom edges, do not match, a direct Fourier
estimate can contain edge artifacts.

The code therefore changes strategy after constructing the binned Fourier maps.
It transforms each binned map back to real space:
\[
  X^{(b)}_{a,c}(x)
  =
  \mathcal{F}^{-1}
  \left[
    M_b(k)\widehat X_{a,c}(k)
  \right](x).
\]

Then, instead of averaging in Fourier space, it forms a real-space product
between the binned first channel and the unbinned second channel:
\[
  X^{(b)}_{a,c_1}(x)\,\overline{X_{a,c_2}(x)}.
\]

\subsection*{Crop mask}

The binned inverse transform $X^{(b)}$ is affected by boundaries. The code
therefore builds a bin-dependent square crop mask
\[
  W_b(x)\in\{0,1\},
\]
which is zero near the border and one in the retained central region.

The non-periodic estimate is then approximately
\[
  P^{\mathrm{nonpbc}}_{a,c_1,c_2,b}
  =
  A_b
  \frac{
    \sum_x
    W_b(x)\,
    X^{(b)}_{a,c_1}(x)\,
    \overline{X_{a,c_2}(x)}
  }{
    \sum_k M_b(k)
  }.
\]
The factor $A_b$ compensates for the fact that only part of the real-space
image is kept:
\[
  A_b
  =
  \frac{N^2}{\sum_x W_b(x)}.
\]

\subsection*{Interpretation}

This branch tries to estimate the same radial cross-power while reducing edge
contamination. It does this by:
\begin{enumerate}[label=\arabic*.]
  \item filtering one channel by radial Fourier bin,
  \item transforming the filtered map back to real space,
  \item multiplying by the second channel in real space,
  \item averaging only over a central square region.
\end{enumerate}

The important difference from \texttt{PBC=True} is the square crop mask
$W_b(x)$. This makes the non-periodic estimate spatially non-uniform: pixels
inside the crop contribute, pixels outside the crop do not.

\section{Why the two branches can behave differently}

With \texttt{PBC=True}, the power-spectrum term is spatially uniform. Every
Fourier mode contributes through the radial mask, and there is no real-space
crop:
\[
  \texttt{PBC=True}
  \quad\Rightarrow\quad
  \text{Fourier average only.}
\]

With \texttt{PBC=False}, the power-spectrum term contains bin-dependent
real-space square masks:
\[
  \texttt{PBC=False}
  \quad\Rightarrow\quad
  \text{Fourier binning} + \text{real-space cropped average.}
\]

Therefore, when the power-spectrum loss is backpropagated to the recovered
maps $u_Q,u_U$, the \texttt{PBC=False} branch can imprint the geometry of the
crop masks into the gradient. This is a plausible source of square or block-like
features in recovered maps when \texttt{ST\_COMPUTE\_PS=True}.

\section{Memory shapes}

Both branches construct binned Fourier maps with shape
\[
  N_b\times C\times B\times N\times N.
\]

They also form channel-pair products with shape
\[
  N_b\times C\times C\times B\times N\times N.
\]

For \texttt{PBC=False}, the code additionally keeps analogous real-space
objects for the inverse-FFT and cropped-product path. This is why the rough
backpropagation memory estimate contains
\[
  2N_bCBN^2 + 2N_bC^2BN^2
\]
complex-valued full-grid quantities.

\end{document}
